% Generated draft: requires mathematical and source review.
\documentclass[11pt]{article}
\usepackage[margin=25mm]{geometry}
\usepackage{fontspec}
\setmainfont{Latin Modern Roman}
\usepackage{amsmath,amssymb,mathtools,unicode-math}
\setmathfont{Latin Modern Math}
\usepackage{enumitem,needspace,xurl,hyperref}
\setlength{\parindent}{0pt}
\setlength{\parskip}{5pt}
\setlength{\emergencystretch}{3em}
\newcommand{\N}{\mathbb{N}}
\newcommand{\Z}{\mathbb{Z}}
\newcommand{\Q}{\mathbb{Q}}
\newcommand{\R}{\mathbb{R}}
\newcommand{\C}{\mathbb{C}}
\newcommand{\F}{\mathbb{F}}
\newcommand{\A}{\mathbb{A}}
\newcommand{\PP}{\mathbb{P}}
\newcommand{\E}{\mathbb{E}}
\newcommand{\eps}{\varepsilon}
\newcommand{\dd}{\,\mathrm{d}}
\begin{document}
\section*{UnsolvedMath ID: 10400157}
\textbf{Source problem code:} catalogue entry
\subsection{Definitions and mathematical statement}
Let \(\Sigma_g\) be a closed oriented surface and let \(\operatorname{Mod}(\Sigma_g)\) be its orientation-preserving mapping class group. In the usual genus-\(g\geq2\) setting, the Nielsen--Thurston classification distinguishes periodic classes (finite order), nonperiodic reducible classes (preserving an essential multicurve), and pseudo-Anosov classes (admitting transverse measured foliations scaled by factors \(\lambda\) and \(\lambda^{-1}\), with \(\lambda>1\)). Low-genus cases require their corresponding conventions.

A TQFT \((V,Z)\) assigns a state space \(V(\Sigma_g)\). Mapping cylinders induce an action, written
\[\rho_{V,Z,g}:\operatorname{Mod}(\Sigma_g)\longrightarrow
\operatorname{PGL}\bigl(V(\Sigma_g)\bigr),\]
when an anomaly makes it projective, or a linear action after the required extra structures or central extension are specified. Masbaum asks whether and how the Nielsen--Thurston type of a mapping class \(f\) relates to its operator \(\rho_{V,Z,g}(f)\). A precise result would state which TQFTs, levels and genera are considered, what operator information is used, and what implications or recovery of mapping-class type it provides.

The excerpt proposes no specific equivalence, spectral criterion or infinite-order assertion. In particular it does not assert that a single TQFT representation distinguishes all three types. An order or spectral statement in the projective setting must be invariant under multiplication of a representative matrix by a nonzero scalar. The quantified family of theories/levels and the meaning of ``relation'' remain open formulation choices in this source request.
\subsection{Short English statement}
Relate the geometric Nielsen--Thurston type of a mapping class to the linear or projective operators supplied by TQFT. The source asks for a relationship without specifying a particular detection theorem or family of theories.
\subsection{Sources}
\begin{itemize}
\item Local UnsolvedMath registry, permanent ID \(10400157\), exact source code \texttt{catalogue entry}.
\item G. Masbaum, attribution in Ohtsuki, \emph{Problems on invariants of knots and 3-manifolds} (2002), Problem 8.11, PDF page 126: \url{https://msp.org/gtm/2002/04/gtm-2002-04-024s.pdf}. Registry access 2026-07-29; pdf-font-structure extraction.
\end{itemize}
\end{document}
